% ---------------------------------------------------------------------------
% 04 — Pitch Deck, 12 slides (ES, beamer). Compilar desde templates/latex:
%   latexmk -pdf 04-pitch-deck/main-es.tex
% Estructura: Sequoia (proposito/problema/solucion/por que ahora/mercado/
% competencia/producto/modelo/equipo/finanzas) + Kawasaki 10/20/30: 10–12
% slides, fuente grande, una idea por slide. Autocontenido: beamer fija su
% propia geometria, por eso NO usa ../common/preambulo (ver RESEARCH.md).
% ---------------------------------------------------------------------------
\documentclass[aspectratio=169,12pt]{beamer}
\usepackage[spanish]{babel}
\usepackage[T1]{fontenc}
\usepackage{lmodern}
\usepackage{booktabs}
\usepackage{graphicx}

% --- Identidad (ver DESIGN.md): crema + coral + tinta ---
\definecolor{bftCoral}{HTML}{CC785C}
\definecolor{bftInk}{HTML}{141413}
\definecolor{bftCanvas}{HTML}{FAF9F5}
\setbeamercolor{background canvas}{bg=bftCanvas}
\setbeamercolor{normal text}{fg=bftInk}
\setbeamercolor{structure}{fg=bftInk}
\setbeamercolor{title}{fg=bftInk}
\setbeamercolor{frametitle}{fg=bftInk}
\setbeamercolor{block title}{bg=bftCoral,fg=white}
\setbeamertemplate{navigation symbols}{}
\setbeamertemplate{footline}[frame number]

\title{[Nombre de la empresa]}
\subtitle{[Propuesta de valor en una frase]}
\author{[Equipo fundador]}
\date{[Mes, anio]}

\begin{document}

% 1/12 — Titulo
\begin{frame}
  \titlepage
\end{frame}

% 2/12 — Proposito
\begin{frame}{Propósito}
  \Large [Declaracion en una frase: que hacemos y para quien.]
\end{frame}

% 3/12 — Problema
\begin{frame}{Problema}
  \begin{itemize}
    \item [Dolor 1 con cifra y fuente]
    \item [Dolor 2: como lo resuelven hoy (hojas de calculo, proveedor legacy)]
  \end{itemize}
\end{frame}

% 4/12 — Solucion
\begin{frame}{Solución}
  [Producto + caso de uso en 2 lineas.] \\[0.5em]
  Diferenciador: [ventaja unica].
\end{frame}

% 5/12 — Por que ahora
\begin{frame}{¿Por qué ahora?}
  \begin{itemize}
    \item [Tendencia 1: cambio tecnologico o regulatorio]
    \item [Tendencia 2: ventana de mercado]
  \end{itemize}
\end{frame}

% 6/12 — Mercado
\begin{frame}{Mercado: TAM / SAM / SOM}
  \begin{center}
  \begin{tabular}{lrr}
    \toprule
    Nivel & Definicion & Valor (MXN) \\
    \midrule
    TAM & Total & [ ] \\
    SAM & Alcanzable & [ ] \\
    SOM & Capturable 3 anios & [ ] \\
    \bottomrule
  \end{tabular}
  \end{center}
\end{frame}

% 7/12 — Competencia
\begin{frame}{Competencia}
  [Rivales clave + nuestras ventajas injustas: IP, efectos de red, costo.]
\end{frame}

% 8/12 — Producto
\begin{frame}{Producto}
  [Funciones clave, arquitectura en una frase, roadmap: MVP → v1 → escala.]
\end{frame}

% 9/12 — Modelo de negocio
\begin{frame}{Modelo de negocio}
  \begin{itemize}
    \item Ingresos: [suscripcion \$49/usuario/mes]
    \item Ticket / LTV: [ ] · CAC: [ ]
    \item Distribucion: [directo / aliados]
  \end{itemize}
\end{frame}

% 10/12 — Traccion y go-to-market
\begin{frame}{Tracción y go-to-market}
  [Usuarios/ARR piloto, crecimiento MoM, alianzas.] Canales: [canal 1, 2].
\end{frame}

% 11/12 — Equipo
\begin{frame}{Equipo}
  \begin{itemize}
    \item [Nombre — rol, ex-(empresa), N anios]
    \item [Nombre — rol, credencial] · Asesores: [ ]
  \end{itemize}
\end{frame}

% 12/12 — Finanzas y solicitud
\begin{frame}{Finanzas y solicitud}
  \begin{itemize}
    \item Ingresos 3 anios: [0.5 / 2 / 5] MM · Margen bruto: [X]\%
    \item Solicitamos: **[monto] MXN** — [40\% producto · 30\% ventas · 30\% operacion]
    \item Contacto: [email / telefono]
  \end{itemize}
\end{frame}

\end{document}
